% TOP_PROBLEM: {"id": 2200008, "title": "Problems Around Polynomials — Problem 4", "category_id": 4, "category": {"id": 4, "name": "algebra", "display_name": "Algebra", "description": "Group theory, ring theory, field theory, and algebraic structures.", "slug": "algebra", "order_index": 4, "created_at": "2026-07-31T15:26:25.671Z"}, "status": "open", "problem_number": "AMR-021-0008", "set_id": 13, "statusQualification": "The source is a localization research problem, not a single yes-or-no conjecture; the necessary relative leading coefficient is made explicit."}
% FIRST_POSTED: 2026-10-01
% ENTRY_KIND: statement_only
% UnsolvedMath ID 2200008; source catalogue code AMR-021-0008
% !TeX program = lualatex
% Definition and sources only; this file is not an LLM solution attempt.
\documentclass[11pt,a4paper]{article}
\usepackage[margin=25mm]{geometry}
\usepackage{fontspec}
\setmainfont{lmroman10-regular.otf}[BoldFont=lmroman10-bold.otf,
 ItalicFont=lmroman10-italic.otf,BoldItalicFont=lmroman10-bolditalic.otf]
\usepackage{amsmath,amssymb,mathtools}
\usepackage{unicode-math}
\setmathfont{latinmodern-math.otf}
\AtBeginDocument{\let\setminus\smallsetminus}
\usepackage{xurl,hyperref}
\hypersetup{colorlinks=true,urlcolor=blue}
\setcounter{secnumdepth}{0}
\setlength{\parindent}{0pt}
\setlength{\parskip}{5pt}
\setlength{\emergencystretch}{3em}
\begin{document}
\section*{Problems Around Polynomials — Problem 4}\label{2200008}
{\small UnsolvedMath ID 2200008 · AMR-021-0008 · Algebra · AMR Open Problem Lists}
\subsection{Definitions and mathematical statement}
Let $p,q\in\mathbb R[z]$ be nonzero polynomials, of arbitrary degrees,
and put $s(z)=p(z)+iq(z)$. Their zero multisets include multiplicities
and are invariant under complex conjugation. The problem asks for
geometric localization and counting results for the complex zeros of
$s$ in terms of these two multisets and the relative leading
coefficient of $p$ and $q$.

The last datum matters: if $P,Q$ are the corresponding monic
polynomials, multiplication of $s$ by a nonzero real constant reduces
the question to
\[
 S_t(z)=P(z)+itQ(z),\qquad t\in\mathbb R\setminus\{0\}.
\]
Thus the input may be specified as two conjugation-invariant finite
multisets and a nonzero real number $t$. Common roots of $P$ and $Q$
are allowed and must be accounted for, with the actual multiplicities
they have in $S_t$.

Find sharp restrictions on which regions of the plane can contain
these zeros and how many zeros can occur in such regions, counted
with multiplicity. Seek conditions stated directly through the root
geometry of $P$ and $Q$, for example their order or separation on the
real line and their placement in half-planes. The classical case of
simple interlacing real zeros motivates the question; the target
allows nonreal and multiple zeros and unequal degrees. A complete
answer should specify both its hypotheses and the region or root
count it guarantees.
\subsection{Short English statement}
Extend root-location rules for a real polynomial plus an imaginary multiple of another real polynomial beyond the interlacing case.
\subsection{Sources}
\begin{itemize}
\item[S1] ULAM.AI and the UnsolvedMath Contributors, supplied problems.json, record 2200008 (AMR-021-0008), AMR Open Problem Lists. Catalogue identity and source transcription; CC BY 4.0.
\url{https://huggingface.co/datasets/ulamai/UnsolvedMath}.
\item[S2] Shapiro - Problems Around Polynomials: The Good, The Bad and The Ugly (2015), source-order item 8 (Problem 4 as printed). Original source indexed by AMR-021-0008.
\url{https://arxiv.org/abs/1503.05295}.
\end{itemize}
\end{document}
